%--- Part 5b: Performance Regimes (6 min) --------------------------------------

\section{Performance Regimes}

\begin{frame}{Analysis Pipeline}
  \begin{center}
    \includegraphics[width=\textwidth]{pipeline-clustering}
  \end{center}
\end{frame}

\renewcommand{\cornerpipeline}{cornerpipeline/pipeline-clustering}

%--- Slide 18 ------------------------------------------------------------------
\begin{frame}{Performance Regime Methodology}
  \begin{enumerate}
    \item Divide trajectories into difficulty groups based on \texttt{ml\_ade}
          (easy / medium / hard)
    \item Run clustering on trajectory + scene characteristics per group
    \begin{itemize}
      \item Option A: original feature space
      \item Option B: UMAP-reduced space
            (dimensions chosen via trustworthiness score + elbow rule)
    \end{itemize}
    \item Assess clustering quality with DBCV score
    \item Export identified clusters (performance regimes)
  \end{enumerate}
\end{frame}

%--- Slide 19 ------------------------------------------------------------------
\begin{frame}{Clustering Results}
  \begin{columns}[t]
    \begin{column}{0.50\textwidth}
      \textbf{Performance group sizes} (VIF-only run)\\[0.3em]
      \renewcommand{\arraystretch}{1.2}
      \begin{tabular}{lccc}
        \toprule
        \textbf{Group} & \textbf{Boundary} & \textbf{N} & \textbf{\%} \\
        \midrule
        Easy   & $<0.39$ & 15\,348 & 59\% \\
        Medium & $0.39$–$1.20$ & 8\,094 & 31\% \\
        Hard   & $>1.20$ & 2\,680 & 10\% \\
        \bottomrule
      \end{tabular}
      \vspace{0.5em}
      \begin{itemize}[<2->]
        \item<2-> Clustering algorithm: \textbf{OPTICS} on raw
              feature-effect space (no UMAP needed for 6-dimensional space)
        \item<2-> Quality assessed via \textbf{DBCV} score
      \end{itemize}
    \end{column}
    \begin{column}{0.47\textwidth}
      \onslide<3->{%
      \textbf{Cluster quality by group}\\[0.3em]
      \renewcommand{\arraystretch}{1.2}
      \begin{tabular}{lcc}
        \toprule
        \textbf{Group} & \textbf{DBCV} & \textbf{Noise \%} \\
        \midrule
        Easy   & \textcolor{lmugreen}{\textbf{0.411}} & 30\% \\
        Medium & 0.397 & 40\% \\
        Hard   & 0.091 & 75\% \\
        \bottomrule
      \end{tabular}
      \vspace{0.4em}
      \footnotesize
      Easy group has the clearest substructure.\\
      Hard is mostly a continuous high-error regime.
      }
    \end{column}
  \end{columns}
\end{frame}

%--- Slide 20 ------------------------------------------------------------------
\begin{frame}{Regime Inspection — What Drives Error in Each \\ Regime?}
  \begin{columns}[t]
    \begin{column}{0.48\textwidth}
      \textbf{Easy regime}
      \begin{itemize}[<+->]
        \item Low \texttt{std\_speed} ($-2.0\sigma$ global) and
              low heading change $\to$ slow, predictable trajectories
        \item Two sub-clusters:
              \emph{slow/stationary} (median ADE 0.02)
              vs.\ \emph{mild motion} (median ADE 0.22)
        \item Both drive negative total feature effects
      \end{itemize}
    \end{column}
    \begin{column}{0.48\textwidth}
      \onslide<2->{\textbf{Hard regime}}
      \begin{itemize}
        \item<2-> 99.5\% of hard trajectories have \emph{positive}
                  total feature effect
        \item<2-> Dominated by high \texttt{std\_speed}
                  ($+2.0\sigma$ global), turning, and acceleration
        \item<3-> No sharp sub-clusters — hard is a
                  \textbf{continuous high-variation regime};
                  a small extreme outlier group has
                  $\texttt{std\_speed} \approx +4.8\sigma$
      \end{itemize}
    \end{column}
  \end{columns}
  \vspace{0.4em}
  \onslide<3->{%
  \footnotesize
  \textbf{Takeaway:} Trajectron++ struggles with erratic, fast-turning pedestrians.
  Slow or near-stationary pedestrians are easy regardless of scene complexity.}
\end{frame}
